\section{光滑结构、流与 Lie 导数}
\begin{align*}
\dd^2&=0,\\
F^*\dd&=\dd F^*,\\
[X,Y](f)&=X(Yf)-Y(Xf),\\
\Lie_Xf&=X(f),\\
\Lie_XY&=[X,Y],\\
\Lie_X\alpha&=\interior{X}\dd\alpha+\dd\interior{X}\alpha.
\end{align*}
主要定义与公式见\eqnrefs{eq:math-dg-exterior-derivative}、\eqnrefs{eq:math-dg-synthesis-lie-bracket}、\eqnrefs{eq:math-dg-synthesis-lie-flow}与\eqnrefs{eq:math-dg-synthesis-cartan}。

\section{标架、联络、挠率与曲率}
\begin{align*}
[e_a,e_b]&=C^c{}_{ab}e_c,\\
\dd\theta^a&=-\frac12C^a{}_{bc}\theta^b\wedge\theta^c,\\
\nabla_{e_b}e_c&=\Gamma^a{}_{bc}e_a,\\
\omega^a{}_c&=\Gamma^a{}_{bc}\theta^b,\\
\Theta^a&=\dd\theta^a+\omega^a{}_b\wedge\theta^b,\\
\Omega^a{}_b&=\dd\omega^a{}_b+\omega^a{}_c\wedge\omega^c{}_b,\\
D\Theta^a&=\Omega^a{}_b\wedge\theta^b,\\
D\Omega^a{}_b&=0.
\end{align*}
若换标架写为 \(e' = e\Lambda\)，则
\[
\omega'=\Lambda^{-1}\omega\Lambda+\Lambda^{-1}\dd\Lambda,
\qquad
\Omega'=\Lambda^{-1}\Omega\Lambda.
\]
对应正文见\eqnrefs{eq:math-dg-frame-structure-equation}、\eqnrefs{eq:math-dg-connection-gauge-transform}、\eqnrefs{eq:math-dg-synthesis-cartan-first}、\eqnrefs{eq:math-dg-synthesis-cartan-second}与\eqnrefs{eq:math-dg-synthesis-bianchi}。

\section{Levi--Civita、Killing 与 Hodge}
Levi--Civita 联络由 Koszul 公式唯一确定：
\[
\begin{aligned}
2g(\nabla_XY,Z)
={}&Xg(Y,Z)+Yg(Z,X)-Zg(X,Y)\\
&+g([X,Y],Z)-g([Y,Z],X)+g([Z,X],Y).
\end{aligned}
\]
坐标中
\[
\Gamma^k{}_{ij}
=\frac12g^{k\ell}
(\partial_i g_{j\ell}+\partial_j g_{i\ell}-\partial_\ell g_{ij}).
\]
Killing、散度与 Hodge 公式为
\begin{align*}
\Lie_Xg&=0,\\
\nabla_iX_j+\nabla_jX_i&=0,\\
\Lie_X\vol_g&=(\divg X)\vol_g,\\
\divg X&=-\delta X^\flat=*\dd*X^\flat,\\
\delta&=(-1)^{n(k+1)+1}*\dd*,\\
\Delta&=\dd\delta+\delta\dd.
\end{align*}
Hodge 分解见\eqnrefs{eq:math-dg-hodge-decomposition}。

\section{Frobenius 与几何演化}
Frobenius 的向量场与余切判据分别为
\[
[\Gamma(\mathcal D),\Gamma(\mathcal D)]\subset\Gamma(\mathcal D),
\qquad
\dd\alpha^a=\eta^a{}_b\wedge\alpha^b.
\]
余维一时可写成
\[
\alpha\wedge\dd\alpha=0.
\]
时间依赖流与度量变分的核心公式为
\begin{align*}
\frac{\dd}{\dd t}(\Phi_{t,s}^*T_t)
&=\Phi_{t,s}^*(\partial_tT_t+\Lie_{V_t}T_t),\\
\partial_tg^{ij}&=-\dot g^{ij},\\
\partial_t\vol_g&=\frac12\operatorname{tr}_g(\dot g)\vol_g,\\
\partial_t\Gamma^k{}_{ij}
&=\frac12g^{k\ell}(\nabla_i \dot g_{j\ell}+\nabla_j \dot g_{i\ell}-\nabla_\ell \dot g_{ij}).
\end{align*}
完整条件见\eqnrefs{eq:math-dg-transport-identity}、\eqnrefs{eq:math-dg-volume-variation}与\eqnrefs{eq:math-dg-connection-variation}。

\section{Gauss--Codazzi--Ricci}
对等距浸入 \(M\hookrightarrow\overline M\)，
\begin{align*}
\overline\nabla_XY&=\nabla_XY+B(X,Y),\\
\overline\nabla_X\xi&=-A_\xi X+\nabla_X^\perp\xi,\\
\overline g(B(X,Y),\xi)&=g(A_\xi X,Y).
\end{align*}
三条兼容方程见\eqnrefs{eq:math-dg-gauss-general}、\eqnrefs{eq:math-dg-codazzi-general}与\eqnrefs{eq:math-dg-ricci-general}。曲面 \(M^2\subset\RR^3\) 的平坦余维一特例为
\[
K=\det A,
\qquad
h_{ij;k}=h_{ik;j}.
\]
